\chapter{Complexity — Leaf}
%%% generated by the blueprint pipeline from TCSlib.Complexity.ClassPSPACE.TQBFARM.Leaf
\section{Overview}

Declarations of the Lean module \texttt{TCSlib.Complexity.ClassPSPACE.TQBFARM.Leaf}.

\section{Declarations}

\begin{lemma}[Reachability after a true input bit]
\lean{Complexity.LogProg.TQBFARM.read_t}
\label{Complexity.LogProg.TQBFARM.read_t}
\leanok
\uses{Complexity.LogProg.AConf, Complexity.LogProg.AReach, Complexity.LogProg.GuardLe, Complexity.LogProg.TQBFARM.A, Complexity.LogProg.TQBFARM.Lb, Complexity.LogProg.TQBFARM.St, Complexity.LogProg.TQBFARM.St.Bd, Complexity.LogProg.TQBFARM.St.toV, Complexity.LogProg.TQBFARM.St.toV_P, Complexity.LogProg.TQBFARM.orc, Complexity.LogProg.TQBFARM.orc_bit, Complexity.LogProg.TQBFARM.rP, Complexity.LogProg.TQBFARM.step_of, Complexity.LogProg.astep, Complexity.LogProg.vword_whole₁}
\difficulty{3}
Let $x$ be a finite bit string, let $B$ be a natural number, let $l,l_1,l_0$ be program
labels, let $st$ be a seven-register state, and let $c$ be an abstract configuration.
Suppose the instruction at $l$ calls the first oracle in whole-input mode with the
input-pointer register as its sole argument, every register of $st$ is at most $B$, and
the bit of $x$ at position $st.P$, interpreted as false when that position is beyond the
end of $x$, is true. If there is a finite run from label $l_1$ with registers
initialized from $st$ to $c$, with every configuration before the endpoint satisfying
the program guard and register bound $B$, then there is such a run from label $l$ to
$c$.
\end{lemma}

\begin{lemma}[Reachability after a false input bit]
\lean{Complexity.LogProg.TQBFARM.read_f}
\label{Complexity.LogProg.TQBFARM.read_f}
\leanok
\uses{Complexity.LogProg.AConf, Complexity.LogProg.AReach, Complexity.LogProg.GuardLe, Complexity.LogProg.TQBFARM.A, Complexity.LogProg.TQBFARM.Lb, Complexity.LogProg.TQBFARM.St, Complexity.LogProg.TQBFARM.St.Bd, Complexity.LogProg.TQBFARM.St.toV, Complexity.LogProg.TQBFARM.St.toV_P, Complexity.LogProg.TQBFARM.orc, Complexity.LogProg.TQBFARM.orc_bit, Complexity.LogProg.TQBFARM.rP, Complexity.LogProg.TQBFARM.step_of, Complexity.LogProg.astep, Complexity.LogProg.vword_whole₁}
\difficulty{3}
Let $x$ be a finite bit string, let $B$ be a natural number, let $l,l_1,l_0$ be program
labels, let $st$ be a seven-register state, and let $c$ be an abstract configuration.
Suppose the instruction at $l$ calls the first oracle in whole-input mode with the
input-pointer register as its sole argument, every register of $st$ is at most $B$, and
the bit of $x$ at position $st.P$, interpreted as false when that position is beyond the
end of $x$, is false. If there is a finite run from label $l_0$ with registers
initialized from $st$ to $c$, with every configuration before the endpoint satisfying
the program guard and register bound $B$, then there is such a run from label $l$ to
$c$.
\end{lemma}

\begin{lemma}[Prepending a successful bit read]
\lean{Complexity.LogProg.TQBFARM.read_b}
\label{Complexity.LogProg.TQBFARM.read_b}
\leanok
\uses{Complexity.LogProg.AConf, Complexity.LogProg.AReach, Complexity.LogProg.GuardLe, Complexity.LogProg.TQBFARM.A, Complexity.LogProg.TQBFARM.Lb, Complexity.LogProg.TQBFARM.St, Complexity.LogProg.TQBFARM.St.Bd, Complexity.LogProg.TQBFARM.St.toV, Complexity.LogProg.TQBFARM.orc, Complexity.LogProg.TQBFARM.rP, Complexity.LogProg.TQBFARM.read_f, Complexity.LogProg.TQBFARM.read_t}
\difficulty{4}
Let $x$ be a finite bit string, let $B$ be a natural number, let $l,l_1,l_0$ be program
labels, let $st$ be a seven-register state, let $c$ be an abstract configuration, and
let $b$ be a Boolean. Suppose the instruction at $l$ calls the first oracle in
whole-input mode with the input-pointer register as its sole argument, every register of
$st$ is at most $B$, and the bit of $x$ at position $st.P$ (taken to be false beyond the
end of $x$) equals $b$. If there is a bounded, precondition-preserving run from label
$l_1$ when $b$ is true, or from label $l_0$ when $b$ is false, with registers
initialized from $st$, to $c$, then there is such a run from label $l$ with those same
initial registers to $c$.
\end{lemma}

\begin{theorem}[Run for counting a unary variable index]
\lean{Complexity.LogProg.TQBFARM.lc_run}
\label{Complexity.LogProg.TQBFARM.lc_run}
\leanok
\uses{Complexity.LogProg.TQBFARM.A, Complexity.LogProg.TQBFARM.Run, Complexity.LogProg.TQBFARM.St, Complexity.LogProg.TQBFARM.St.Bd, Complexity.LogProg.TQBFARM.fin_of, Complexity.LogProg.TQBFARM.read_b, Complexity.LogProg.TQBFARM.read_f, Complexity.LogProg.TQBFARM.read_t, Complexity.LogProg.TQBFARM.step_of, Complexity.LogProg.astep, Complexity.LogProg.reach_of_eq, Complexity.drop_succ_of_drop, Complexity.getD_of_drop, Complexity.lt_of_drop}
\difficulty{5}
Let $x$ be a finite bit string, $B,n$ natural numbers, $s,b$ Booleans, $rest$ a finite
bit string, and $st$ a machine state. Suppose the suffix of $x$ beginning at position
$st.P$ is $n$ true bits followed by a false bit, then $b$, then $rest$; every register
of $st$ is at most $B$; $|x|\le B$; and $st.V\le st.P$. Then the machine has a guarded
run from label $\text{lc}(s)$ and state $st$ to label $\text{lv}(s,b)$ and the state
obtained by setting $P=st.P+n+2$ and $V=st.V+n$, leaving the other registers unchanged.
\end{theorem}

\begin{theorem}[Run through literal evaluation]
\lean{Complexity.LogProg.TQBFARM.lv_run}
\label{Complexity.LogProg.TQBFARM.lv_run}
\leanok
\uses{Complexity.LogProg.AReach.trans, Complexity.LogProg.TQBFARM.A, Complexity.LogProg.TQBFARM.Run, Complexity.LogProg.TQBFARM.St, Complexity.LogProg.TQBFARM.St.Bd, Complexity.LogProg.TQBFARM.asg, Complexity.LogProg.TQBFARM.cpDst, Complexity.LogProg.TQBFARM.cpExit, Complexity.LogProg.TQBFARM.cpSrc, Complexity.LogProg.TQBFARM.cp_run, Complexity.LogProg.TQBFARM.fin_of, Complexity.LogProg.TQBFARM.hv_run, Complexity.LogProg.TQBFARM.sb_run, Complexity.LogProg.TQBFARM.step_of, Complexity.LogProg.astep}
\difficulty{6}
Let $x$ be a bit string, $B\in\mathbb{N}$, $b\in\{0,1\}$, and let $st=(K,S,P,T,Z,C,V)$
be a seven-register state. If every register is at most $B$ and $T=Z=C=0$, the machine
has a guarded run from label $lv(0,b)$ with state $st$ to label
$cl(\mathrm{asg}(K,S,V)=b)$ with state $(K,S,P,T,Z,C,0)$.
\end{theorem}

\begin{theorem}[Run through a serialized literal]
\lean{Complexity.LogProg.TQBFARM.lit_run}
\label{Complexity.LogProg.TQBFARM.lit_run}
\leanok
\uses{Complexity.LogProg.AReach.trans, Complexity.LogProg.TQBFARM.A, Complexity.LogProg.TQBFARM.Run, Complexity.LogProg.TQBFARM.St, Complexity.LogProg.TQBFARM.St.Bd, Complexity.LogProg.TQBFARM.asg, Complexity.LogProg.TQBFARM.fin_of, Complexity.LogProg.TQBFARM.lc_run, Complexity.LogProg.TQBFARM.lv_run, Complexity.LogProg.TQBFARM.step_of, Complexity.LogProg.astep, Complexity.LogProg.reach_of_eq, Complexity.drop_succ_of_drop, Complexity.le_of_drop, Std.Sat.CNF.serializeLit}
\difficulty{5}
Let $x$ be a bit string, $B,v\in\mathbb{N}$, $s,b\in\{0,1\}$, and let
$st=(K,S,P,T,Z,C,V)$ be a seven-register state. Suppose the suffix of $x$ starting at
position $P$ is the serialization $1^{v+1}0b$ followed by a bit string $rest$, every
register of $st$ is at most $B$, and $|x|\leq B$. If the scratch registers $T,Z,C,V$ are
zero, the machine has a guarded run from label $lt(s)$ with state $st$ to label
$cl(s\lor(\mathrm{asg}(K,S,v)=b))$ with state $(K,S,P+v+3,T,Z,C,V)$.
\end{theorem}

\begin{theorem}[Clause evaluation during a run]
\lean{Complexity.LogProg.TQBFARM.clause_run}
\label{Complexity.LogProg.TQBFARM.clause_run}
\leanok
\uses{Complexity.LogProg.AReach.trans, Complexity.LogProg.TQBFARM.Run, Complexity.LogProg.TQBFARM.St, Complexity.LogProg.TQBFARM.St.Bd, Complexity.LogProg.TQBFARM.asg, Complexity.LogProg.TQBFARM.fin_of, Complexity.LogProg.TQBFARM.lit_run, Complexity.LogProg.TQBFARM.read_f, Complexity.LogProg.TQBFARM.read_t, Complexity.LogProg.reach_of_eq, Complexity.drop_add_of_drop, Complexity.getD_of_drop, Complexity.le_of_drop, Std.Sat.CNF.serializeLit}
\difficulty{5}
Let $x$ be a Boolean word, $B$ a natural number, $rest$ a Boolean word, $Cl$ a clause
over natural-number variables, $s$ a Boolean value, and $st$ a machine state. Suppose
the suffix of $x$ beginning at position $st.P$ is the serialized literals of $Cl$,
followed by false and then $rest$. Suppose every register of $st$ is at most $B$,
$|x|\le B$, and $st.T=st.Z=st.C=st.V=0$. The machine can reach state $ce\bigl(s\lor
\operatorname{eval}(Cl,\operatorname{asg}(st.K,st.S))\bigr)$ from state $cl(s)$, with
all registers unchanged except that the input pointer becomes $st.P+|Cl\text{'s
serialized literals}|$.
\end{theorem}

\begin{theorem}[Evaluation of a Serialized CNF Matrix]
\lean{Complexity.LogProg.TQBFARM.formula_run}
\label{Complexity.LogProg.TQBFARM.formula_run}
\leanok
\uses{Complexity.LogProg.AReach.trans, Complexity.LogProg.TQBFARM.A, Complexity.LogProg.TQBFARM.Run, Complexity.LogProg.TQBFARM.St, Complexity.LogProg.TQBFARM.St.Bd, Complexity.LogProg.TQBFARM.asg, Complexity.LogProg.TQBFARM.clause_run, Complexity.LogProg.TQBFARM.fin_of, Complexity.LogProg.TQBFARM.read_f, Complexity.LogProg.TQBFARM.read_t, Complexity.LogProg.TQBFARM.step_of, Complexity.LogProg.astep, Complexity.LogProg.reach_of_eq, Complexity.drop_add_of_drop, Complexity.drop_succ_of_drop, Complexity.getD_of_drop, Complexity.le_of_drop, Std.Sat.CNF.serialize, Std.Sat.CNF.serializeClause, Std.Sat.CNF.serializeLit}
\difficulty{5}
Let $x$ be a Boolean input word, $B$ a natural number, $r$ a suffix, $varphi$ a CNF
formula over natural-number variables, and $	ext{st}$ a machine state. Suppose that
deleting the first $	ext{st}.P$ bits of $x$ leaves the serialization of $varphi$
followed by $r$, each of the seven registers of $	ext{st}$ is at most $B$, the length of
$x$ is at most $B$, and $	ext{st}.T=	ext{st}.Z=	ext{st}.C=	ext{st}.V=0$. Then the
machine can reach the return label with value $varphi$ evaluated under the assignment
determined by $	ext{st}.K$ and $	ext{st}.S$, in a configuration obtained from $	ext{st}$
by setting only $P$ to zero. Every configuration before the endpoint satisfies the
machine guard and has every register at most $B$.
\end{theorem}
